\scanpage{85}

quand la suite $(a_{\sigma i})$ est décroissante, on est donc \deleted{\unsure{en changeant} au besoin \unsure{de notations}} \deleted{\unsure{supposer, pour prouver (8), que la suite} $a_1\geq a_2\mathord{-}\geq a_n$} ramené à prouver que si $a_1\geq a_2\mathord{-}\geq a_n$, on a $\sum a_i\alpha_i\leq\deleted{\unsure{a_i\alpha|a_i\alpha_i}}\sum a_i\beta_i\deleted{\sum a_i\unsure{\alpha|\alpha_i}}$, \deleted{\unsure{en supposant} \unsure{de plus}}\inserted{pourvu que} les $a_i$ soient positifs \deleted{si on n'est pas dans le cas b) du lemme} ou qu'on ait $\sum\alpha_i=\sum\beta_i$. Or on a en effet

\[
\begin{array}{@{}l@{\qquad}l@{}}
(9)&\displaystyle\sum a_i\unsure{\alpha|\beta}_i=(a_1-a_2)\unsure{\alpha|\beta}_1+(a_2-a_3)(\unsure{\alpha|\beta}_1+\unsure{\alpha|\beta}_2)+\mathord{-}+(a_{n-1}-a_n)(\unsure{\alpha|\beta}_1+\mathord{-}+\unsure{\alpha|\beta}_{n-1})\\
&\displaystyle\hfill+a_n(\unsure{\alpha|\beta}_1+\mathord{-}+\unsure{\alpha|\beta}_n)
\end{array}
\]

\deleted{\unsure{et on} $\leq$} et une formule analogue pour $\deleted{\sum a_i\unsure{\alpha_i}}\sum a_i\beta_i$. Comme les coefficients $(a_1-a_2)\mathord{.}\,\mathord{-}\,\mathord{-}(a_i-a_{i-1}),\mathord{-}\,\mathord{-}(a_{n-1}-a_n)$ sont positifs, \unsure{on} \deleted{\unsure{on conclut que la somme relative}} et que $\deleted{\sum}\alpha_1\leq\beta_1,\mathord{.}\,\mathord{-}\,\mathord{.}(\alpha_1+\mathord{-}+\alpha_{n-1})\leq(\beta_1+\mathord{-}+\beta_{n-1})$, on en conclut que tous les termes de la somme relative aux $\alpha_i$ sont majorés par les termes corres-\newline pondants de la somme relative aux $\beta_i$, à l'exception \inserted{à priori} du dernier. Mais on a aussi $a_n(\alpha_1\mathord{-}\alpha_n)\leq a_n(\beta_1+\mathord{-}\unsure{+}\beta_{\unsure{n|i}})$ dans \deleted{\unsure{du terme du dernier}}\inserted{le cas a)} du lemme ($a_n\geq0$, $\sum\alpha_i\leq\sum\beta$ et dans le cas b) ($\sum\alpha_i=\sum\beta_i$), d'où \deleted{la conclusion}\inserted{l'inégalité (1)}.

\deleted{\unsure{Notons}} Signalons en passant le

{\reversemarginpar\margin{\uline{petits caractères}}}
{\reversemarginpar\margin{\begin{tikzpicture}[x=.009cm,y=-.009cm]
\draw[line width=.35pt] (46,29) .. controls (56,42) and (54,88) .. (53,133) .. controls (52,192) and (47,231) .. (35,266) .. controls (23,303) and (21,350) .. (26,396) .. controls (30,448) and (45,496) .. (50,545) .. controls (56,594) and (46,639) .. (39,684) .. controls (28,741) and (26,781) .. (29,817) .. controls (31,837) and (38,852) .. (42,865);
\draw[line width=.45pt] (121,19) .. controls (122,58) and (109,104) .. (111,153) .. controls (112,176) and (115,203) .. (111,227);
\end{tikzpicture}}}

\uline{Corollaire \unsure{1a|bis|1bis}} Sous les conditions de \deleted{la prop. \unsure{du} \illegible{1}, b,} \deleted{\illegible{1}} \deleted{\unsure{corollaire 1}} \inserted{\uline{Th. 1, 2°, b)}} \deleted{\unsure{(resp. du corollaire 1, a))}} \deleted{la prop. \unsure{corollaire}, b)} on aura aussi, \deleted{\unsure{si on a pour tout} $\unsure{\sigma}$}

\[
\begin{array}{@{}l@{\qquad}l@{}}
(10)&\displaystyle f\left(\frac1{\sigma_1},\mathord{-}\,\mathord{-}\frac1{\sigma_n}\right)\leq f\left(\frac1{\rho_1},\mathord{-}\,\mathord{-},\frac1{\rho_n}\right)
\end{array}
\]

En effet, de $\rho_1\mathord{-}\rho_n=\sigma_1\mathord{-}\sigma_n$ on tire pour tout
{\reversemarginpar\margin{\inserted{$1\leq M\leq n$ :}}}

\[
\rho_M\rho_{M+1}\mathord{-}\rho_n=\frac{\sigma_1\mathord{-}\sigma_{M-1}}{\rho_1\mathord{-}\rho_{M-1}}\sigma_{\unsure{M}}\mathord{-}\,\mathord{-}\sigma_n,\qquad\text{d'où }\rho_M\mathord{-}\rho_n\leq\sigma_M\mathord{-}\,\mathord{-}\sigma_n
\]

en vertu de $\sigma_1\mathord{-}\sigma_{M-1}\leq\rho_1\mathord{-}\,\mathord{-}\rho_{M-1}$. Il s'ensuit que que les suites $(\frac1{\rho_n},\mathord{-}\frac1{\rho_1})$ et $((\frac1{\unsure{\rho_n|\sigma_n}},\mathord{.}\,\mathord{-}\frac1{\sigma_1})$ satisfont encore aux conditions du corollaire (la première étant décroissante, les produits \deleted{\unsure{qui interviennent par}} des $\frac1{\rho_i}$ ég. au produit des $\frac1{\sigma_i}$, et enfin $\frac1{\rho_n}\mathord{-}\,\mathord{-}\frac1{\rho_{\unsure{i}}}\leq\frac1{\sigma_{\unsure{n}}}\mathord{-}\,\mathord{-}\frac1{\sigma_{\unsure{i}}}$ pour $1\leq M\leq n$), d'où l'inégalité \unsure{(10)}.

{\reversemarginpar\margin{86}}
